\documentclass[10pt,a4paper]{amsart}
\usepackage{amsmath,amssymb,amsthm,graphicx}
\usepackage{verbatim}
\newenvironment{reference}{\comment}{\endcomment}
%\newenvironment{reference}{}{}
\newenvironment{slogan}{\comment}{\endcomment}
\newenvironment{history}{\comment}{\endcomment}

% For commutative diagrams we use Xy-pic
\usepackage[all]{xy}

% We use 2cell for 2-commutative diagrams.
\xyoption{2cell}
% Theorem environments.
%
\theoremstyle{plain}
\newtheorem{theorem}[subsection]{Theorem}
\newtheorem{proposition}[subsection]{Proposition}
\newtheorem{lemma}[subsection]{Lemma}

\theoremstyle{definition}
\newtheorem{definition}[subsection]{Definition}
\newtheorem{example}[subsection]{Example}
\newtheorem{exercise}[subsection]{Exercise}
\newtheorem{situation}[subsection]{Situation}

\theoremstyle{remark}
\newtheorem{remark}[subsection]{Remark}
\newtheorem{remarks}[subsection]{Remarks}

\numberwithin{equation}{subsection}

% Macros
%
\def\lim{\mathop{\mathrm{lim}}\nolimits}
\def\colim{\mathop{\mathrm{colim}}\nolimits}
\def\Spec{\mathop{\mathrm{Spec}}}
\def\Hom{\mathop{\mathrm{Hom}}\nolimits}
\def\Ext{\mathop{\mathrm{Ext}}\nolimits}
\def\SheafHom{\mathop{\mathcal{H}\!\mathit{om}}\nolimits}
\def\SheafExt{\mathop{\mathcal{E}\!\mathit{xt}}\nolimits}
\def\Sch{\mathit{Sch}}
\def\Mor{\mathop{\mathrm{Mor}}\nolimits}
\def\Ob{\mathop{\mathrm{Ob}}\nolimits}
\def\Sh{\mathop{\mathit{Sh}}\nolimits}
\def\NL{\mathop{N\!L}\nolimits}
\def\CH{\mathop{\mathrm{CH}}\nolimits}
\def\proetale{{pro\text{-}\acute{e}tale}}
\def\etale{{\acute{e}tale}}
\def\QCoh{\mathit{QCoh}}
\def\Ker{\mathop{\mathrm{Ker}}}
\def\Im{\mathop{\mathrm{Im}}}
\def\Coker{\mathop{\mathrm{Coker}}}
\def\Coim{\mathop{\mathrm{Coim}}}

% Boxtimes
%
\DeclareMathSymbol{\boxtimes}{\mathbin}{AMSa}{"02}

%
% Macros for moduli stacks/spaces
%
\def\QCohstack{\mathcal{QC}\!\mathit{oh}}
\def\Cohstack{\mathcal{C}\!\mathit{oh}}
\def\Spacesstack{\mathcal{S}\!\mathit{paces}}
\def\Quotfunctor{\mathrm{Quot}}
\def\Hilbfunctor{\mathrm{Hilb}}
\def\Curvesstack{\mathcal{C}\!\mathit{urves}}
\def\Polarizedstack{\mathcal{P}\!\mathit{olarized}}
\def\Complexesstack{\mathcal{C}\!\mathit{omplexes}}
% \Pic is the operator that assigns to X its picard group, usage \Pic(X)
% \Picardstack_{X/B} denotes the Picard stack of X over B
% \Picardfunctor_{X/B} denotes the Picard functor of X over B
\def\Pic{\mathop{\mathrm{Pic}}\nolimits}
\def\Picardstack{\mathcal{P}\!\mathit{ic}}
\def\Picardfunctor{\mathrm{Pic}}
\def\Deformationcategory{\mathcal{D}\!\mathit{ef}}
\usepackage{cleveref}
\crefname{originalsection}{original section}{original sections}
\Crefname{originalsection}{Original section}{Original sections}
\makeatletter
\newlabel{section-intrinsic}{{03DG}{source}{}{sourcepointer-1}{}}
\newlabel{section-intrinsic@cref}{{[originalsection][0][]03DG}{[1][0][]source}{}{}{}}
\makeatother
\crefname{originaldefinition}{original definition}{original definitions}
\Crefname{originaldefinition}{Original definition}{Original definitions}
\makeatletter
\newlabel{definition-site-local}{{03DL}{source}{}{sourcepointer-2}{}}
\newlabel{definition-site-local@cref}{{[originaldefinition][0][]03DL}{[1][0][]source}{}{}{}}
\makeatother
\crefname{originalreference}{original reference}{original references}
\Crefname{originalreference}{Original reference}{Original references}
\makeatletter
\newlabel{sites-lemma-localize-morphism-strong}{{03EF}{source}{}{sourcepointer-3}{}}
\newlabel{sites-lemma-localize-morphism-strong@cref}{{[originalreference][0][]03EF}{[1][0][]source}{}{}{}}
\makeatother
\crefname{originallemma}{original lemma}{original lemmas}
\Crefname{originallemma}{Original lemma}{Original lemmas}
\makeatletter
\newlabel{lemma-global-pullback}{{03DF}{source}{}{sourcepointer-4}{}}
\newlabel{lemma-global-pullback@cref}{{[originallemma][0][]03DF}{[1][0][]source}{}{}{}}
\makeatother
\crefname{originallemma}{original lemma}{original lemmas}
\Crefname{originallemma}{Original lemma}{Original lemmas}
\makeatletter
\newlabel{lemma-local-final-object}{{03DN}{source}{}{sourcepointer-5}{}}
\newlabel{lemma-local-final-object@cref}{{[originallemma][0][]03DN}{[1][0][]source}{}{}{}}
\makeatother
\title{Pullback along a morphism of ringed topoi preserves the local properties of modules}
\author{The Stacks Project Authors}
\date{stacks-project:sites-modules.tex:local-snapshot-20261012; GNU FDL 1.2 or later}
\begin{document}
\maketitle
\noindent\textit{Editorial note.} All delivered mathematical text is the author's own, copied verbatim from the pinned Stacks Project chapter file sites-modules.tex; no statement or proof step was rewritten, reordered or removed.\par\medskip
\noindent\textit{Editorial note.} The counted span cites five external items (sections and lemmas of the same chapter family), recorded in assembly.external\_references with canonical tags as navigation pointers.\par\medskip
\noindent\textit{Allowed standard prerequisites (editorial).}\begin{itemize}
\item Definitions of the local properties of modules on a site. Stacks Project, Modules on Sites, Definition definition-site-local (tag 09B9). Lists the seven local properties whose preservation is proved, and their global counterparts.
\item The section on intrinsic properties, reducing ringed-topoi statements to left exact continuous functors. Stacks Project, Modules on Sites, Section section-intrinsic (tag 09BB). Justifies the initial reduction to morphisms of ringed sites with a left exact continuous underlying functor.
\end{itemize}\medskip
\ifdefined\hypertarget\hypertarget{sourcepointer-1}{}\fi
\noindent\textit{Source reference.} Original section 03DG: sources/stacks-chow/tags.txt, source lines 4419--4419. Cross-chapter standard prerequisite cited by the author inside the counted span; it is used as a pointer and its own proof is not part of this item.\par\medskip
\ifdefined\hypertarget\hypertarget{sourcepointer-2}{}\fi
\noindent\textit{Source reference.} Original definition 03DL: sources/stacks-chow/tags.txt, source lines 4424--4424. Cross-chapter standard prerequisite cited by the author inside the counted span; it is used as a pointer and its own proof is not part of this item.\par\medskip
\ifdefined\hypertarget\hypertarget{sourcepointer-3}{}\fi
\noindent\textit{Source reference.} Original reference 03EF: sources/stacks-chow/tags.txt, source lines 4454--4454. Cross-chapter standard prerequisite cited by the author inside the counted span; it is used as a pointer and its own proof is not part of this item.\par\medskip
\ifdefined\hypertarget\hypertarget{sourcepointer-4}{}\fi
\noindent\textit{Source reference.} Original lemma 03DF: sources/stacks-chow/tags.txt, source lines 4418--4418. Cross-chapter standard prerequisite cited by the author inside the counted span; it is used as a pointer and its own proof is not part of this item.\par\medskip
\ifdefined\hypertarget\hypertarget{sourcepointer-5}{}\fi
\noindent\textit{Source reference.} Original lemma 03DN: sources/stacks-chow/tags.txt, source lines 4426--4426. Cross-chapter standard prerequisite cited by the author inside the counted span; it is used as a pointer and its own proof is not part of this item.\par\medskip
\begin{lemma}
\label{lemma-local-pullback}
Let
$(f, f^\sharp) :
(\Sh(\mathcal{C}), \mathcal{O}_\mathcal{C})
\to
(\Sh(\mathcal{D}), \mathcal{O}_\mathcal{D})$
be a morphism of ringed topoi.
Let $\mathcal{F}$ be an $\mathcal{O}_\mathcal{D}$-module.
\begin{enumerate}
\item If $\mathcal{F}$ is locally free then $f^*\mathcal{F}$ is locally free.
\item If $\mathcal{F}$ is finite locally free then $f^*\mathcal{F}$ is
finite locally free.
\item If $\mathcal{F}$ is locally generated by sections
then $f^*\mathcal{F}$ is locally generated by sections.
\item If $\mathcal{F}$ is locally generated by $r$ sections
then $f^*\mathcal{F}$ is locally generated by $r$ sections.
\item If $\mathcal{F}$ is of finite type
then $f^*\mathcal{F}$ is of finite type.
\item If $\mathcal{F}$ is quasi-coherent then
$f^*\mathcal{F}$ is quasi-coherent.
\item If $\mathcal{F}$ is of finite presentation
then $f^*\mathcal{F}$ is of finite presentation.
\end{enumerate}
\end{lemma}

\begin{proof}
According to the discussion in Section \ref{section-intrinsic}
we need only check preservation under pullback for a morphism of ringed sites
$(f, f^\sharp) :
(\mathcal{C}, \mathcal{O}_\mathcal{C})
\to
(\mathcal{D}, \mathcal{O}_\mathcal{D})$
such that $f$ is given by a left exact, continuous functor
$u : \mathcal{D} \to \mathcal{C}$ between sites which have
all finite limits.
Let $\mathcal{G}$ be a sheaf of $\mathcal{O}_\mathcal{D}$-modules
which has one of the properties (1) -- (6) of
Definition \ref{definition-site-local}.
We know $\mathcal{D}$ has a final object $Y$ and $X = u(Y)$
is a final object for $\mathcal{C}$. By assumption we have
a covering $\{Y_i \to Y\}$ such that $\mathcal{G}|_{\mathcal{D}/Y_i}$
has the corresponding global property. Set $X_i = u(Y_i)$ so
that $\{X_i \to X\}$ is a covering in $\mathcal{C}$.
We get a commutative diagram of morphisms ringed sites
$$
\xymatrix{
(\mathcal{C}/X_i, \mathcal{O}_\mathcal{C}|_{X_i}) \ar[r] \ar[d] &
(\mathcal{C}, \mathcal{O}_\mathcal{C}) \ar[d] \\
(\mathcal{D}/Y_i, \mathcal{O}_\mathcal{D}|_{Y_i}) \ar[r] &
(\mathcal{D}, \mathcal{O}_\mathcal{D})
}
$$
by Sites, Lemma \ref{sites-lemma-localize-morphism-strong}.
Hence by Lemma \ref{lemma-global-pullback}
that $f^*\mathcal{G}|_{X_i}$ has the corresponding global
property. Hence we conclude that $\mathcal{G}$ has the local
property we started out with by Lemma \ref{lemma-local-final-object}.
\end{proof}

\end{document}
